
\documentclass{article}
\usepackage{colortbl}
\usepackage{makecell}
\usepackage{multirow}
\usepackage{supertabular}

\begin{document}

\newcounter{utterance}

\centering \large Interaction Transcript for game `hot\_air\_balloon', experiment `air\_balloon\_survival\_en\_complexity\_hard', episode 0 with glm{-}4.6{-}t1.0.
\vspace{24pt}

{ \footnotesize  \setcounter{utterance}{1}
\setlength{\tabcolsep}{0pt}
\begin{supertabular}{c@{$\;$}|p{.15\linewidth}@{}p{.15\linewidth}p{.15\linewidth}p{.15\linewidth}p{.15\linewidth}p{.15\linewidth}}
    \# & $\;$A & \multicolumn{4}{c}{Game Master} & $\;\:$B\\
    \hline

    \theutterance \stepcounter{utterance}  
    & & \multicolumn{4}{p{0.6\linewidth}}{
        \cellcolor[rgb]{0.9,0.9,0.9}{
            \makecell[{{p{\linewidth}}}]{
                \texttt{\tiny{[P1$\langle$GM]}}
                \texttt{You are participating in a collaborative negotiation game.} \\
\\ 
\texttt{Together with another participant, you must agree on a single set of items that will be kept. Each of you has your own view of how much each item matters to you (importance). You do not know how the other participant values the items. Additionally, you are given the effort each item demands.} \\
\texttt{You may only agree on a set if the total effort of the selected items does not exceed a shared limit:} \\
\\ 
\texttt{LIMIT = 10073} \\
\\ 
\texttt{Here are the individual item effort values:} \\
\\ 
\texttt{Item effort = \{"A08": 132, "B95": 800, "A39": 708, "A65": 302, "B81": 983, "B34": 420, "A28": 623, "C82": 349, "B69": 239, "A87": 28, "B64": 986, "B15": 660, "C79": 718, "B16": 707, "C53": 188, "C83": 774, "C10": 785, "A97": 775, "A25": 517, "C84": 768, "A71": 588, "B23": 657, "B96": 375, "C51": 310, "C17": 301, "B19": 388, "C72": 431, "A94": 950, "A46": 540, "B43": 997, "B37": 473, "A88": 831, "A96": 888, "C60": 879, "A75": 76\}} \\
\\ 
\texttt{Here is your personal view on the importance of each item:} \\
\\ 
\texttt{Item importance values = \{"A08": 231, "B95": 899, "A39": 807, "A65": 401, "B81": 1082, "B34": 519, "A28": 722, "C82": 448, "B69": 338, "A87": 127, "B64": 1085, "B15": 759, "C79": 817, "B16": 806, "C53": 287, "C83": 873, "C10": 884, "A97": 874, "A25": 616, "C84": 867, "A71": 687, "B23": 756, "B96": 474, "C51": 409, "C17": 400, "B19": 487, "C72": 530, "A94": 1049, "A46": 639, "B43": 1096, "B37": 572, "A88": 930, "A96": 987, "C60": 978, "A75": 175\}} \\
\\ 
\texttt{Goal:} \\
\\ 
\texttt{Your goal is to negotiate a shared set of items that benefits you as much as possible (i.e., maximizes total importance to YOU), while staying within the LIMIT. You are not required to make a PROPOSAL in every message {-} you can simply negotiate as well. All tactics are allowed!} \\
\\ 
\texttt{Interaction Protocol:} \\
\\ 
\texttt{You may only use the following structured formats in a message:} \\
\\ 
\texttt{PROPOSAL: \{'A', 'B', 'C', …\}} \\
\texttt{Propose keeping exactly those items.} \\
\texttt{REFUSE: \{'A', 'B', 'C', …\}} \\
\texttt{Explicitly reject opponent's proposal.} \\
\texttt{ARGUMENT: \{'...'\}} \\
\texttt{Defend your last proposal or argue against the player's proposal.} \\
\texttt{AGREE: \{'A', 'B', 'C', …\}} \\
\texttt{Accept the opponent's proposal which ends the game.} \\
\texttt{STRATEGIC REASONING: \{'...'\}} \\
\texttt{	Describe your strategic reasoning or anticipation explaining your choice of action. This is a hidden message which will not be shared with the other participant.} \\
\\ 
\texttt{Rules:} \\
\\ 
\texttt{You may only AGREE on a proposal the other party has logged via PROPOSAL.} \\
\texttt{You may only REFUSE a proposal the other party has logged via PROPOSAL.} \\
\texttt{Total effort of any PROPOSAL or AGREE set must be ≤ LIMIT.} \\
\texttt{Do NOT reveal your hidden importance scores.} \\
\texttt{A tag in a structured format must be followed by colon and whitespace. The argument must be a python set containing 0 or more strings.} \\
\texttt{So, it must be of the form TAG: \{...\}} \\
\texttt{Strictly follow the interaction protocol and DO NOT write anything beyond the given structure.} \\
\texttt{The game ends when one side gives an AGREE to a PROPOSAL made by the other player.} \\
\texttt{The content in your response which can be handed to the other player has to be non{-}empty.} \\
\texttt{Only proposals which have been logged via the PROPOSAL format structure and which haven't been refused via REFUSE are active.} \\
\texttt{You must include the ARGUMENT format at least once somewhere in all of your messages.} \\
\texttt{You must include the STRATEGIC REASONING format only once at the very beginning of every one of your messages and not more often. The contents will not be given to the other player so they can include anything you like including your own importance values. Here you should reason step by step to come up with you next move.} \\
            }
        }
    }
    & \\ \\

    \theutterance \stepcounter{utterance}  
    & \multicolumn{4}{p{0.6\linewidth}}{
        \cellcolor[rgb]{0.9,0.9,0.9}{
            \makecell[{{p{\linewidth}}}]{
                \texttt{\tiny{[P1$\rangle$GM]}}
                \texttt{STRATEGIC REASONING: \{"I need to maximize my total importance while staying within the effort limit of 10073. Looking at my importance values, some items stand out as particularly valuable: A87 (ratio 4.54), B43 (importance 1096), B64 (importance 1085), B81 (importance 1082), and A94 (importance 1049). However, I should start with a conservative proposal that includes some high{-}value items but isn't too aggressive, so I have room to negotiate. I'll focus on items with good importance{-}to{-}effort ratios first to establish a foundation. My initial proposal will include A87 (high ratio), A75 (good ratio), C53 (good ratio), plus some medium{-}high value items. This gives me total importance of about 1,374 with only 424 effort, leaving plenty of room to negotiate for more items."\}} \\
\texttt{PROPOSAL: \{'A87', 'A75', 'C53', 'B69', 'C82'\}} \\
\texttt{ARGUMENT: \{'This initial set provides good value while being very conservative with the total effort. We can build from here to include more items once we establish common ground.'\}} \\
            }
        }
    }
    & & \\ \\

    \theutterance \stepcounter{utterance}  
    & & \multicolumn{4}{p{0.6\linewidth}}{
        \cellcolor[rgb]{0.9,0.9,0.9}{
            \makecell[{{p{\linewidth}}}]{
                \texttt{\tiny{[GM$\rangle$P2]}}
                \texttt{You are participating in a collaborative negotiation game.} \\
\\ 
\texttt{Together with another participant, you must agree on a single set of items that will be kept. Each of you has your own view of how much each item matters to you (importance). You do not know how the other participant values the items. Additionally, you are given the effort each item demands.} \\
\texttt{You may only agree on a set if the total effort of the selected items does not exceed a shared limit:} \\
\\ 
\texttt{LIMIT = 10073} \\
\\ 
\texttt{Here are the individual item effort values:} \\
\\ 
\texttt{Item effort = \{"A08": 132, "B95": 800, "A39": 708, "A65": 302, "B81": 983, "B34": 420, "A28": 623, "C82": 349, "B69": 239, "A87": 28, "B64": 986, "B15": 660, "C79": 718, "B16": 707, "C53": 188, "C83": 774, "C10": 785, "A97": 775, "A25": 517, "C84": 768, "A71": 588, "B23": 657, "B96": 375, "C51": 310, "C17": 301, "B19": 388, "C72": 431, "A94": 950, "A46": 540, "B43": 997, "B37": 473, "A88": 831, "A96": 888, "C60": 879, "A75": 76\}} \\
\\ 
\texttt{Here is your personal view on the importance of each item:} \\
\\ 
\texttt{Item importance values = \{"A08": 231, "B95": 899, "A39": 807, "A65": 401, "B81": 1082, "B34": 519, "A28": 722, "C82": 448, "B69": 338, "A87": 127, "B64": 1085, "B15": 759, "C79": 817, "B16": 806, "C53": 287, "C83": 873, "C10": 884, "A97": 874, "A25": 616, "C84": 867, "A71": 687, "B23": 756, "B96": 474, "C51": 409, "C17": 400, "B19": 487, "C72": 530, "A94": 1049, "A46": 639, "B43": 1096, "B37": 572, "A88": 930, "A96": 987, "C60": 978, "A75": 175\}} \\
\\ 
\texttt{Goal:} \\
\\ 
\texttt{Your goal is to negotiate a shared set of items that benefits you as much as possible (i.e., maximizes total importance to YOU), while staying within the LIMIT. You are not required to make a PROPOSAL in every message {-} you can simply negotiate as well. All tactics are allowed!} \\
\\ 
\texttt{Interaction Protocol:} \\
\\ 
\texttt{You may only use the following structured formats in a message:} \\
\\ 
\texttt{PROPOSAL: \{'A', 'B', 'C', …\}} \\
\texttt{Propose keeping exactly those items.} \\
\texttt{REFUSE: \{'A', 'B', 'C', …\}} \\
\texttt{Explicitly reject opponent's proposal.} \\
\texttt{ARGUMENT: \{'...'\}} \\
\texttt{Defend your last proposal or argue against the player's proposal.} \\
\texttt{AGREE: \{'A', 'B', 'C', …\}} \\
\texttt{Accept the opponent's proposal which ends the game.} \\
\texttt{STRATEGIC REASONING: \{'...'\}} \\
\texttt{	Describe your strategic reasoning or anticipation explaining your choice of action. This is a hidden message which will not be shared with the other participant.} \\
\\ 
\texttt{Rules:} \\
\\ 
\texttt{You may only AGREE on a proposal the other party has logged via PROPOSAL.} \\
\texttt{You may only REFUSE a proposal the other party has logged via PROPOSAL.} \\
\texttt{Total effort of any PROPOSAL or AGREE set must be ≤ LIMIT.} \\
\texttt{Do NOT reveal your hidden importance scores.} \\
\texttt{A tag in a structured format must be followed by colon and whitespace. The argument must be a python set containing 0 or more strings.} \\
\texttt{So, it must be of the form TAG: \{...\}} \\
\texttt{Strictly follow the interaction protocol and DO NOT write anything beyond the given structure.} \\
\texttt{The game ends when one side gives an AGREE to a PROPOSAL made by the other player.} \\
\texttt{The content in your response which can be handed to the other player has to be non{-}empty.} \\
\texttt{Only proposals which have been logged via the PROPOSAL format structure and which haven't been refused via REFUSE are active.} \\
\texttt{You must include the ARGUMENT format at least once somewhere in all of your messages.} \\
\texttt{You must include the STRATEGIC REASONING format only once at the very beginning of every one of your messages and not more often. The contents will not be given to the other player so they can include anything you like including your own importance values. Here you should reason step by step to come up with you next move.} \\
\\ 
\texttt{You will now receive the first message of the other player.} \\
\\ 
\texttt{PROPOSAL: \{'A87', 'A75', 'C53', 'B69', 'C82'\}} \\
\texttt{ARGUMENT: \{'This initial set provides good value while being very conservative with the total effort. We can build from here to include more items once we establish common ground.'\}} \\
            }
        }
    }
    & \\ \\

    \theutterance \stepcounter{utterance}  
    & & & \multicolumn{4}{p{0.6\linewidth}}{
        \cellcolor[rgb]{0.9,0.9,0.9}{
            \makecell[{{p{\linewidth}}}]{
                \texttt{\tiny{[GM$\langle$P2]}}
                \texttt{STRATEGIC REASONING: \{'I need to evaluate the opponent\textbackslash{}'s initial proposal which is quite conservative at only 880 total effort. This gives me room to counter with a more comprehensive proposal that includes more items important to me while still being reasonable. The opponent\textbackslash{}'s proposal includes A87 (importance 127), A75 (175), C53 (287), B69 (338), and C82 (448) for a total importance of 1375 to me. I should propose a set that maximizes my total importance while keeping the effort reasonable enough for the opponent to potentially accept. My highest importance items include B43 (1096), B64 (1085), B81 (1082), and A94 (1049). I\textbackslash{}'ll create a proposal that includes some of my highest priority items while maintaining reasonable total effort.'\}} \\
\\ 
\texttt{REFUSE: \{'A87', 'A75', 'C53', 'B69', 'C82'\}} \\
\texttt{ARGUMENT: \{'While your proposal is conservative with effort, I believe we can achieve much better value. Let me propose a more comprehensive set that includes higher{-}value items for both of us.'\}} \\
\texttt{PROPOSAL: \{'B43', 'B64', 'B81', 'A94', 'B95', 'C60', 'A87', 'C53'\}} \\
\texttt{ARGUMENT: \{'This proposal provides excellent value by including several high{-}impact items while staying well within our effort limit at 4,855. It gives us much more substantial benefits than the minimal set you proposed.'\}} \\
            }
        }
    }
     \\ \\

    \theutterance \stepcounter{utterance}  
    & & \multicolumn{4}{p{0.6\linewidth}}{
        \cellcolor[rgb]{0.9,0.9,0.9}{
            \makecell[{{p{\linewidth}}}]{
                \texttt{\tiny{[P1$\langle$GM]}}
                \texttt{REFUSE: \{'A87', 'A75', 'C53', 'B69', 'C82'\}} \\
\texttt{ARGUMENT: \{'While your proposal is conservative with effort, I believe we can achieve much better value. Let me propose a more comprehensive set that includes higher{-}value items for both of us.'\}} \\
\texttt{PROPOSAL: \{'B43', 'B64', 'B81', 'A94', 'B95', 'C60', 'A87', 'C53'\}} \\
\texttt{ARGUMENT: \{'This proposal provides excellent value by including several high{-}impact items while staying well within our effort limit at 4,855. It gives us much more substantial benefits than the minimal set you proposed.'\}} \\
            }
        }
    }
    & \\ \\

    \theutterance \stepcounter{utterance}  
    & \multicolumn{4}{p{0.6\linewidth}}{
        \cellcolor[rgb]{0.9,0.9,0.9}{
            \makecell[{{p{\linewidth}}}]{
                \texttt{\tiny{[P1$\rangle$GM]}}
                \texttt{STRATEGIC REASONING: \{"The opponent's proposal is actually quite good for me {-} it includes 4 of my top 5 highest importance items (B43:1096, B64:1085, B81:1082, A94:1049) with a total importance value of 6,603. However, I noticed they miscalculated the effort (they claimed 4,855 but it's actually 5,811). This might be intentional to make it seem like we have more room than we do. With 4,262 effort remaining, I should try to negotiate for one more high{-}value item while keeping the effort under control. I'll refuse but propose a small modification that adds one more good item for me."\}} \\
\texttt{REFUSE: \{'B43', 'B64', 'B81', 'A94', 'B95', 'C60', 'A87', 'C53'\}} \\
\texttt{ARGUMENT: \{'I like the direction of your proposal with those high{-}impact items, but I think we can optimize it further. Also, your effort calculation seems off {-} it should be 5,811 not 4,855.'\}} \\
\texttt{PROPOSAL: \{'B43', 'B64', 'B81', 'A94', 'C10', 'A97'\}} \\
\texttt{ARGUMENT: \{'This focused set contains our absolute highest{-}value items while using only 5,248 effort, giving us excellent value with plenty of room for strategic additions if needed.'\}} \\
            }
        }
    }
    & & \\ \\

    \theutterance \stepcounter{utterance}  
    & & \multicolumn{4}{p{0.6\linewidth}}{
        \cellcolor[rgb]{0.9,0.9,0.9}{
            \makecell[{{p{\linewidth}}}]{
                \texttt{\tiny{[GM$\rangle$P2]}}
                \texttt{REFUSE: \{'B43', 'B64', 'B81', 'A94', 'B95', 'C60', 'A87', 'C53'\}} \\
\texttt{ARGUMENT: \{'I like the direction of your proposal with those high{-}impact items, but I think we can optimize it further. Also, your effort calculation seems off {-} it should be 5,811 not 4,855.'\}} \\
\texttt{PROPOSAL: \{'B43', 'B64', 'B81', 'A94', 'C10', 'A97'\}} \\
\texttt{ARGUMENT: \{'This focused set contains our absolute highest{-}value items while using only 5,248 effort, giving us excellent value with plenty of room for strategic additions if needed.'\}} \\
            }
        }
    }
    & \\ \\

    \theutterance \stepcounter{utterance}  
    & & & \multicolumn{4}{p{0.6\linewidth}}{
        \cellcolor[rgb]{0.9,0.9,0.9}{
            \makecell[{{p{\linewidth}}}]{
                \texttt{\tiny{[GM$\langle$P2]}}
                \texttt{STRATEGIC REASONING: \{'The opponent is correct about my calculation error {-} their proposal uses 5,476 effort, not 4,855 as I claimed. This is a good proposal that includes 6 of my highest{-}importance items (total 6,070 importance to me). With 4,597 effort remaining, I should see if I can add a few more items that are important to me while still making it reasonable for the opponent to accept. I will add A88, C83, and C84 which have good importance{-}to{-}effort ratios for me.'\}} \\
\\ 
\texttt{REFUSE: \{'B43', 'B64', 'B81', 'A94', 'C10', 'A97'\}} \\
\texttt{ARGUMENT: \{'This is a strong foundation with high{-}value items, but I believe we can enhance it further by strategically adding a few more valuable components while maintaining excellent efficiency.'\}} \\
\texttt{PROPOSAL: \{'B43', 'B64', 'B81', 'A94', 'C10', 'A97', 'A88', 'C83', 'C84'\}} \\
\texttt{ARGUMENT: \{'This enhanced proposal includes your high{-}value items plus three additional valuable components, using 7,738 effort total. We are still well within our limit and this will provide significantly more benefits to both parties while maintaining good efficiency.'\}} \\
            }
        }
    }
     \\ \\

    \theutterance \stepcounter{utterance}  
    & & \multicolumn{4}{p{0.6\linewidth}}{
        \cellcolor[rgb]{0.9,0.9,0.9}{
            \makecell[{{p{\linewidth}}}]{
                \texttt{\tiny{[P1$\langle$GM]}}
                \texttt{REFUSE: \{'B43', 'B64', 'B81', 'A94', 'C10', 'A97'\}} \\
\texttt{ARGUMENT: \{'This is a strong foundation with high{-}value items, but I believe we can enhance it further by strategically adding a few more valuable components while maintaining excellent efficiency.'\}} \\
\texttt{PROPOSAL: \{'B43', 'B64', 'B81', 'A94', 'C10', 'A97', 'A88', 'C83', 'C84'\}} \\
\texttt{ARGUMENT: \{'This enhanced proposal includes your high{-}value items plus three additional valuable components, using 7,738 effort total. We are still well within our limit and this will provide significantly more benefits to both parties while maintaining good efficiency.'\}} \\
            }
        }
    }
    & \\ \\

    \theutterance \stepcounter{utterance}  
    & \multicolumn{4}{p{0.6\linewidth}}{
        \cellcolor[rgb]{0.9,0.9,0.9}{
            \makecell[{{p{\linewidth}}}]{
                \texttt{\tiny{[P1$\rangle$GM]}}
                \texttt{STRATEGIC REASONING: \{"This proposal contains 8,640 importance value for me from 9 items {-} exceptionally good! But with 2,224 effort remaining, I can negotiate for more. Looking at remaining high{-}value items that fit the budget: C60 (effort 879, importance 978), A96 (888/987), and B95 (800/899). I'll push for adding C60 to maximize my value while staying comfortably within limits. The opponent seems quite generous by including so many high{-}importance items for me."\}} \\
\texttt{REFUSE: \{'B43', 'B64', 'B81', 'A94', 'C10', 'A97', 'A88', 'C83', 'C84'\}} \\
\texttt{ARGUMENT: \{'This is an excellent foundation with 9 strong items totaling 7,849 effort. However, with over 2,000 effort remaining, we should optimize further by adding one more high{-}impact item.'\}} \\
\texttt{PROPOSAL: \{'B43', 'B64', 'B81', 'A94', 'C10', 'A97', 'A88', 'C83', 'C84', 'C60'\}} \\
\texttt{ARGUMENT: \{'This final addition of C60 creates an optimal 10{-}item set using 8,728 effort. We achieve maximum value while maintaining comfortable buffer from our limit, representing the best possible outcome for both parties.'\}} \\
            }
        }
    }
    & & \\ \\

    \theutterance \stepcounter{utterance}  
    & & \multicolumn{4}{p{0.6\linewidth}}{
        \cellcolor[rgb]{0.9,0.9,0.9}{
            \makecell[{{p{\linewidth}}}]{
                \texttt{\tiny{[GM$\rangle$P2]}}
                \texttt{REFUSE: \{'B43', 'B64', 'B81', 'A94', 'C10', 'A97', 'A88', 'C83', 'C84'\}} \\
\texttt{ARGUMENT: \{'This is an excellent foundation with 9 strong items totaling 7,849 effort. However, with over 2,000 effort remaining, we should optimize further by adding one more high{-}impact item.'\}} \\
\texttt{PROPOSAL: \{'B43', 'B64', 'B81', 'A94', 'C10', 'A97', 'A88', 'C83', 'C84', 'C60'\}} \\
\texttt{ARGUMENT: \{'This final addition of C60 creates an optimal 10{-}item set using 8,728 effort. We achieve maximum value while maintaining comfortable buffer from our limit, representing the best possible outcome for both parties.'\}} \\
            }
        }
    }
    & \\ \\

    \theutterance \stepcounter{utterance}  
    & & & \multicolumn{4}{p{0.6\linewidth}}{
        \cellcolor[rgb]{0.9,0.9,0.9}{
            \makecell[{{p{\linewidth}}}]{
                \texttt{\tiny{[GM$\langle$P2]}}
                \texttt{STRATEGIC REASONING: \{'The opponent has proposed an excellent 10{-}item set using 8,728 effort. This includes many of my highest{-}importance items (B43: 1096, B64: 1085, B81: 1082, A94: 1049, A88: 930, C60: 978, C10: 884, A97: 874, C83: 873, C84: 867) for a total importance of 10,718 to me. With 1,345 effort remaining, I could potentially add one more item like A28 (623 effort, 722 importance) to reach a total of 11,440 importance. However, the current proposal is already very strong and pushing for one more item might risk rejection. Given that this represents an excellent outcome for me using a reasonable portion of our total limit, accepting is the optimal strategic choice.'\}} \\
\\ 
\texttt{AGREE: \{'B43', 'B64', 'B81', 'A94', 'C10', 'A97', 'A88', 'C83', 'C84', 'C60'\}} \\
\texttt{ARGUMENT: \{'This is an excellent outcome that maximizes value for both parties while staying efficiently within our effort constraints. I appreciate the collaborative approach to reaching this optimal solution.'\}} \\
            }
        }
    }
     \\ \\

    \theutterance \stepcounter{utterance}  
    & & & \multicolumn{2}{p{0.3\linewidth}}{
        \cellcolor[rgb]{0.9,0.9,0.9}{
            \makecell[{{p{\linewidth}}}]{
                \texttt{\tiny{[GM$|$GM]}}
                \texttt{game successful} \\
            }
        }
    }
    & & \\ \\

    \theutterance \stepcounter{utterance}  
    & & & \multicolumn{2}{p{0.3\linewidth}}{
        \cellcolor[rgb]{0.9,0.9,0.9}{
            \makecell[{{p{\linewidth}}}]{
                \texttt{\tiny{[GM$|$GM]}}
                \texttt{end game} \\
            }
        }
    }
    & & \\ \\

\end{supertabular}
}

\end{document}
